\section{Introdução}

Este trabalho propõe um aplicativo móvel integrado a \textit{smartwatches} que transforma uma corrida em uma experiência de sobrevivência. O usuário tenta fugir de uma horda virtual de zumbis, e a posição do personagem no jogo é determinada pela distância percorrida, pelo ritmo de corrida e por sinais biométricos. O objetivo do aplicativo é aumentar o engajamento em atividades físicas, reduzir a sensação de monotonia e promover a adesão a longo prazo por meio de feedback em tempo real e desafios adaptados.

A proposta utiliza uma narrativa imersiva em que o jogador percebe-se perseguido pela horda de zumbis, acompanhando o avanço do progresso e a proximidade dos inimigos. Dessa forma, a gamificação deixa de ser apenas um elemento estético e se integra à mecânica principal da corrida.

A prática regular de atividade física é amplamente reconhecida como um pilar fundamental para a saúde e o bem-estar. Ainda assim, manter o engajamento e garantir a adesão ao longo do tempo continua a ser um desafio significativo. As altas taxas de desistência da prática de exercícios são em grande parte motivadas pela chamada "barreira do tédio", o que faz com que atividades habituais passem a ser vistas como monótonas, repetitivas e sem recompensas positivas imediatas. Na fase inicial de qualquer mudança de comportamento, a motivação intrínseca é tipicamente baixa. Consequentemente, a carência de um retorno interativo propicia o abandono prematuro \cite{plangger2022little}. \\ 
\indent Para diminuir esse distanciamento, a gamificação, isto é, a aplicação de mecânicas de design de jogos em contextos que não envolvem jogos, tem se consolidado como um recurso poderoso para promover mudanças de comportamento na saúde. Nesse cenário emergiram os exergames (jogos de exercício), que utilizam a tecnologia para transformar o esforço físico em experiências interativas e estimulantes \cite{li2025gamified}. Apesar de promissoras, muitas abordagens de gamificação genéricas enfrentam limitações atreladas ao "efeito de novidade" (onde o interesse tende a diminuir após as primeiras semanas) por usarem dificuldades estáticas, que não se adequam às necessidades individuais do usuário \cite{li2025gamified, helmi2024mediating}. \\
\indent A fim de contornar essas limitações e criar um engajamento sustentável, a literatura atual aponta para a integração de tecnologias vestíveis, como \textit{smartwatches}, que possibilitam o monitoramento fisiológico contínuo e a aplicação do \textit{biofeedback} \cite{bastos2023gamification, shukurova2025ai}. O acesso a esses dados biométricos em tempo real permite a implementação de algoritmos de Ajuste Dinâmico de Dificuldade (\textit{DDA} — \textit{Dynamic Difficulty Adjustment}), um mecanismo inteligente que altera os desafios do jogo de forma responsiva ao nível de habilidade e esforço cardiovascular atual do praticante \cite{li2025gamified}. \\
\indent Diante deste contexto, este artigo propõe o desenvolvimento de uma aplicação móvel integrada a \textit{smartwatches} com foco na gamificação de atividades físicas. O sistema utiliza o \textit{biofeedback} em tempo real para estruturar uma narrativa de sobrevivência imersiva, na qual o desempenho biológico e o movimento do usuário ditam diretamente a progressão no jogo. A arquitetura também prevê mecanismos de adaptação de dificuldade, utilizados nesta etapa de forma controlada a partir dos dados de telemetria coletados, com o objetivo de aproximar a experiência do usuário do estado de fluxo, além de promover o condicionamento seguro e superar a barreira do tédio, consolidando a adesão a longo prazo.
